% GENERATED FILE. Edit canonical lessons, then run npm run generate:books.
% canonical-source-hash: fnv1a64:cefa266f32d6992b
% canonical-lessons: ES-C672-indicador, ES-C672-letrero, ES-C672-flecha, ES-C672-senalizacion, ES-C672-desvio

\chapter{Carteles y flechas}
\label{ch:es-carteles-y-flechas}
\hlchaptermodality{672}

\begin{chapteropening}
\textbf{By the end of this chapter:} \emph{I can read signs on the way: signposts, notices, arrows, the signposting and diversions.}

Complete the last lesson of chapter 672: I can read signs on the way: signposts, notices, arrows, the signposting and diversions.
\end{chapteropening}

\section[el indicador]{el indicador — the signpost, the sign; the gauge}

\label{lesson:ES-C672-indicador}



\begin{quote}
Before the new one: say the Spanish for forward, then the Spanish for the detour.
\end{quote}

\subsection*{You'll want to know: el indicador}
\textbf{el indicador} — \textquotedblleft{}the signpost, the sign; the gauge\textquotedblright{}.

\begin{quote}
Sigue los indicadores hacia el centro.
\end{quote}

\begin{quote}
\emph{Follow the signs for the centre.}
\end{quote}

In a car, \emph{el indicador} also shows the fuel or the speed.

\begin{cousinweb}
\emph{indicador} is built on \emph{indicar} \textquotedblleft{}to point out\textquotedblright{}, a learned borrowing from Latin \emph{indicāre} (Lewis \& Short).
\end{cousinweb}

\begin{grammarlens}[title={What you've built}]
One more word for telling, explaining and describing at B1.
\end{grammarlens}

\subsection*{Your turn}
\begin{itemize}
\raggedright
  \item \emph{Say it:} \emph{el indicador}
  \item \emph{Say it:} \emph{el indicador}, once more
  \item \emph{Say it:} a sentence of your own with \emph{el indicador}
  \item \emph{Recall:} say the Spanish for forward, then the Spanish for the detour, then say \emph{el indicador} again
\end{itemize}

\begin{tcolorbox}[breakable,colback=teal!4,colframe=teal!35!black,title={Before you move on}]
What does el indicador mean? (\textquotedblleft{}The signpost, the sign; the gauge\textquotedblright{}.) Say it once more.
\end{tcolorbox}
\section[el letrero]{el letrero — the sign (with words), the notice}

\label{lesson:ES-C672-letrero}



\begin{quote}
Before the new one: say the Spanish for the detour, then the Spanish for the signpost.
\end{quote}

\subsection*{You'll want to know: el letrero}
\textbf{el letrero} — \textquotedblleft{}the sign (with words), the notice\textquotedblright{}.

\begin{quote}
El letrero de la tienda es azul.
\end{quote}

\begin{quote}
\emph{The shop sign is blue.}
\end{quote}

Close to \emph{el cartel}; a \emph{letrero} is the one made of letters.

\begin{cousinweb}
\emph{letrero} is built on \emph{letra} \textquotedblleft{}letter\textquotedblright{} with \emph{-ero} (EtymDB), and \emph{letra} is inherited from Latin \emph{littera}.
\end{cousinweb}

\begin{grammarlens}[title={What you've built}]
One more word for telling, explaining and describing at B1.
\end{grammarlens}

\subsection*{Your turn}
\begin{itemize}
\raggedright
  \item \emph{Say it:} \emph{el letrero}
  \item \emph{Say it:} \emph{el letrero}, once more
  \item \emph{Say it:} a sentence of your own with \emph{el letrero}
  \item \emph{Recall:} say the Spanish for the detour, then the Spanish for the signpost, then say \emph{el letrero} again
\end{itemize}

\begin{tcolorbox}[breakable,colback=teal!4,colframe=teal!35!black,title={Before you move on}]
What does el letrero mean? (\textquotedblleft{}The sign (with words), the notice\textquotedblright{}.) Say it once more.
\end{tcolorbox}
\section[la flecha]{la flecha — the arrow}

\label{lesson:ES-C672-flecha}



\begin{quote}
Before the new one: say the Spanish for the signpost, then the Spanish for the sign.
\end{quote}

\subsection*{You'll want to know: la flecha}
\textbf{la flecha} — \textquotedblleft{}the arrow\textquotedblright{}.

\begin{quote}
Sigue las flechas amarillas.
\end{quote}

\begin{quote}
\emph{Follow the yellow arrows.}
\end{quote}

On a map or a sign, the arrow shows the way.

\begin{cousinweb}
\emph{flecha} is usually traced to French \emph{flèche} \textquotedblleft{}arrow\textquotedblright{}, a word of Germanic origin.
\end{cousinweb}

\begin{grammarlens}[title={What you've built}]
One more word for telling, explaining and describing at B1.
\end{grammarlens}

\subsection*{Your turn}
\begin{itemize}
\raggedright
  \item \emph{Say it:} \emph{la flecha}
  \item \emph{Say it:} \emph{la flecha}, once more
  \item \emph{Say it:} a sentence of your own with \emph{la flecha}
  \item \emph{Recall:} say the Spanish for the signpost, then the Spanish for the sign, then say \emph{la flecha} again
\end{itemize}

\begin{tcolorbox}[breakable,colback=teal!4,colframe=teal!35!black,title={Before you move on}]
What does la flecha mean? (\textquotedblleft{}The arrow\textquotedblright{}.) Say it once more.
\end{tcolorbox}
\section[la señalización]{la señalización — the signs, the signposting}

\label{lesson:ES-C672-senalizacion}



\begin{quote}
Before the new one: say the Spanish for the sign, then the Spanish for the arrow.
\end{quote}

\subsection*{You'll want to know: la señalización}
\textbf{la señalización} — \textquotedblleft{}the signs, the signposting\textquotedblright{}.

\begin{quote}
La señalización de la ciudad es muy clara.
\end{quote}

\begin{quote}
\emph{The signposting in the city is very clear.}
\end{quote}

All the road signs together.

\begin{cousinweb}
\emph{señalización} is built on \emph{señalizar} (EtymDB), from \emph{señal} \textquotedblleft{}sign\textquotedblright{}, inherited from Late Latin \emph{signāle}, from \emph{signum} \textquotedblleft{}a mark\textquotedblright{}.
\end{cousinweb}

\begin{grammarlens}[title={What you've built}]
One more word for telling, explaining and describing at B1.
\end{grammarlens}

\subsection*{Your turn}
\begin{itemize}
\raggedright
  \item \emph{Say it:} \emph{la señalización}
  \item \emph{Say it:} \emph{la señalización}, once more
  \item \emph{Say it:} a sentence of your own with \emph{la señalización}
  \item \emph{Recall:} say the Spanish for the sign, then the Spanish for the arrow, then say \emph{la señalización} again
\end{itemize}

\begin{tcolorbox}[breakable,colback=teal!4,colframe=teal!35!black,title={Before you move on}]
What does la señalización mean? (\textquotedblleft{}The signs, the signposting\textquotedblright{}.) Say it once more.
\end{tcolorbox}
\section[el desvío]{el desvío — the diversion, the detour}

\label{lesson:ES-C672-desvio}



\begin{quote}
Before the new one: say the Spanish for the arrow, then the Spanish for the signs.
\end{quote}

\subsection*{You'll want to know: el desvío}
\textbf{el desvío} — \textquotedblleft{}the diversion, the detour\textquotedblright{}.

\begin{quote}
Hay un desvío por obras.
\end{quote}

\begin{quote}
\emph{There is a diversion because of roadworks.}
\end{quote}

The verb is \emph{desviar}, to divert.

\begin{cousinweb}
\emph{desvío} is built on \emph{desviar} \textquotedblleft{}to divert\textquotedblright{}, from Late Latin \emph{dēviāre} \textquotedblleft{}to leave the road\textquotedblright{}, built on \emph{via} \textquotedblleft{}road\textquotedblright{}.
\end{cousinweb}

\begin{grammarlens}[title={What you've built}]
One more word for telling, explaining and describing at B1.
\end{grammarlens}

\subsection*{Your turn}
\begin{itemize}
\raggedright
  \item \emph{Say it:} \emph{el desvío}
  \item \emph{Say it:} \emph{el desvío}, once more
  \item \emph{Say it:} a sentence of your own with \emph{el desvío}
  \item \emph{Recall:} say the Spanish for the arrow, then the Spanish for the signs, then say \emph{el desvío} again
\end{itemize}

\begin{tcolorbox}[breakable,colback=teal!4,colframe=teal!35!black,title={Before you move on}]
What does el desvío mean? (\textquotedblleft{}The diversion, the detour\textquotedblright{}.) Say it once more.
\end{tcolorbox}
